%% annex/p01.tex --- proofs retained from the verified Part-I substrate.
%% Phase-4's P1 work order explicitly requires an annex home for CSB and
%% Cantor; the other items preserve every migrated Part-I complete proof.
\hbannexchapter{1}{Notation, Symbols, and Conventions}\label{annex:p01}

\begin{annexitem}{annex:p1.1}{thm:p01-ordered-pair-identity}{Ordered-pair identity}
\begin{proof}
The reverse implication follows by substitution.  Suppose
\(\langle a,b\rangle=\langle c,d\rangle\), using the registered
Kuratowski coding.  If \(a=b\), the left side is the singleton
\(\{\{a\}\}\), so the right side is a singleton; hence \(c=d\), and
\(\{\{c\}\}=\{\{a\}\}\) gives \(c=d=a=b\).  If \(a\ne b\), the left
side has the unique singleton member \(\{a\}\) and the distinct two-element
member \(\{a,b\}\).  The right side cannot be a singleton, so \(c\ne d\);
matching its unique singleton member gives \(a=c\), and then
\(\{a,b\}=\{a,d\}\), with \(b\ne a\), gives \(b=d\).
\end{proof}
\end{annexitem}

\begin{annexitem}{annex:p1.2}{thm:p01-csb}{Cantor--Schr\"oder--Bernstein theorem}
\begin{proof}
Let \(A_0=A\setminus g[B]\), let \(A_{n+1}=g[f[A_n]]\), and put
\(C=\bigcup_{n\in\mathbb N}A_n\).  Define
\[
 h(a)=
 \begin{cases}
   f(a),&a\in C,\\
   g^{-1}(a),&a\notin C.
 \end{cases}
\]
The second branch is defined: \(A_0\subseteq C\), so \(a\notin C\)
implies \(a\in g[B]\), and injectivity of \(g\) makes its preimage unique.
Each branch is injective.  Their ranges are disjoint, because
\(a\in C\Rightarrow g(f(a))\in C\), whereas a value \(b=g^{-1}(a')\)
from the second branch satisfies \(g(b)=a'\notin C\).

For surjectivity, take \(b\in B\).  If \(b\in f[C]\), it lies in the first
range.  Otherwise \(g(b)\notin C\): it is not in \(A_0\), since it lies in
\(g[B]\), and membership in any
\(A_{n+1}=g[f[A_n]]\) would, by injectivity of \(g\), imply
\(b\in f[C]\).  Thus \(h(g(b))=b\), so \(h\) is surjective and hence
bijective.
\end{proof}
\end{annexitem}

\begin{annexitem}{annex:p1.3}{thm:p01-cantor}{Cantor's theorem}
\begin{proof}
For any \(f:A\to\mathcal P(A)\), let
\(D=\{a\in A:a\notin f(a)\}\).  If \(f\) were surjective, some
\(d\in A\) would satisfy \(f(d)=D\), whence
\[
 d\in D\quad\Longleftrightarrow\quad d\notin f(d)
 \quad\Longleftrightarrow\quad d\notin D,
\]
impossible.  Thus no surjection exists.  The singleton map
\(a\mapsto\{a\}\) injects \(A\) into \(\mathcal P(A)\), so the cardinal
inequality is strict.
\end{proof}
\end{annexitem}

\begin{annexitem}{annex:p1.4}{thm:p01-syntax-counting}{Syntax counting}
\begin{proof}
Fix an injection \(e:A\to\kappa\).  Applying it coordinatewise gives a
length-preserving injection \(A^{<\omega}\to\kappa^{<\omega}\); the cited
finite-coding theorem gives
\(\abs{\kappa^{<\omega}}=\kappa\), proving the first assertion.

Finite strings over \(\Sigma\) lie in \(\Sigma^{<\omega}\), and finite
tuples of strings lie in \((\Sigma^{<\omega})^{<\omega}\), so the first
assertion applies twice.  Let \(X\) be the at most countable set of logical
symbols, variables, and punctuation adjoined by the fixed grammar.  Since
\(\kappa\) is infinite, item~4 of the cited finite-coding theorem applied to
\(\Sigma,X,\varnothing,\ldots\) gives
\(\abs{\Sigma\cup X}\leq\kappa\).  Every expression generated by a
finitary grammar is a finite string over \(\Sigma\cup X\), so the first
assertion yields the claimed bound.
\end{proof}
\end{annexitem}

\begin{annexitem}{annex:p1.5}{thm:p01-countable-strings}{Countability of finite strings}
\begin{proof}
First verify the pairing map used below.  Put \(t_s=s(s+1)/2\).  For every
\(z\in\mathbb N\), the consecutive finite intervals
\([t_s,t_{s+1})\) give a unique \(s\) with
\(t_s\leq z<t_{s+1}\).  Set \(b=z-t_s\) and \(a=s-b\).  The upper
inequality gives \(0\leq b\leq s\), so \(a,b\in\mathbb N\), and direct
substitution gives \(\operatorname{pair}_{\mathbb N}(a,b)=z\).
Conversely, if \(z=\operatorname{pair}_{\mathbb N}(a,b)\), then
\(s=a+b\) is the unique interval index, and the displayed inverse recovers
\(b\) and then \(a\).  Thus the map is bijective.

Now fix an injection \(e:A\to\mathbb N\).  Define injections
\(e_n:A^n\to\mathbb N\) recursively by
\(e_0(\epsilon)=0\) and
\[
 e_{n+1}(w)=\operatorname{pair}_{\mathbb N}
   \bigl(e_n(w\mathbin{\upharpoonright}n),e(w(n))\bigr).
\]
Induction on \(n\) proves injectivity: equality of paired values gives
equality of the last symbols and, through \(e_n\), of the initial segments.
Finally send \(w\in A^n\) to
\(\operatorname{pair}_{\mathbb N}(n,e_n(w))\).  Equality of outputs recovers
the length and then the string, so this is an injection \(A^*\to\mathbb N\).
\end{proof}
\end{annexitem}

\begin{annexitem}{annex:p1.6}{thm:p01-countable-formulas}{Countability of formulas over a countable alphabet}
\begin{proof}
Every generated formula is a finite string over the alphabet.  The formula
set is therefore a subset of the countable set \(A^*\) from
Lemma~\ref{thm:p01-countable-strings}; restricting its injection into
\(\mathbb N\) proves the claim.
\end{proof}
\end{annexitem}

\begin{annexitem}{annex:p1.7}{thm:p01-concatenation-laws}{Concatenation laws}
\begin{proof}
Let \(m,n,p\in\mathbb N\) be the unique lengths for which
\(u\in A^m\), \(v\in A^n\), and \(w\in A^p\).  Both
\((uv)w\) and \(u(vw)\) have domain \(m+n+p\).  At an index
\(i<m+n+p\), the three exhaustive cases are
\(i<m\), \(m\leq i<m+n\), and \(m+n\leq i<m+n+p\).  The defining clauses
for concatenation give, respectively, the common values
\[
 u(i),\qquad v(i-m),\qquad w(i-(m+n)).
\]
Thus the strings agree pointwise.  For \(\epsilon u\), every index
\(i<m\) enters the second concatenation clause and has value \(u(i-0)\);
for \(u\epsilon\), every such index enters the first clause and has value
\(u(i)\).  Their lengths are \(m\), so both identity equations follow.
\end{proof}
\end{annexitem}

\begin{annexitem}{annex:p1.8}{thm:p01-structural-induction}{Structural induction}
\begin{proof}
Let \(X=\{x\in G:P(x)\}\).  The base hypotheses give \(B\subseteq X\),
and the constructor hypotheses say that \(X\) is closed under every member
of \(F\).  Since \(G\) is the least subset containing \(B\) and closed under
\(F\), \(G\subseteq X\).  The reverse inclusion holds by definition of
\(X\), so every member of \(G\) has property \(P\).
\end{proof}
\end{annexitem}

\begin{annexitem}{annex:p1.9}{thm:p01-structural-recursion}{Structural recursion}
\begin{proof}
Let \(G_0\) be the base elements and let \(G_{m+1}\) contain \(G_m\) and
every constructor value on arguments from \(G_m\).  Generation gives
\(G=\bigcup_mG_m\).  Assign each \(x\in G\) its least stage.  Recursion on
that stage sets \(h(b)\) to the specified base value and, for the unique
construction \(x=f(x_1,\ldots,x_k)\), sets
\[
 h(x)=H_f(h(x_1),\ldots,h(x_k));
\]
the immediate arguments have smaller stage.  Unique construction prevents
incompatible defining clauses, so this gives a total function satisfying the
equations.  If \(h'\) satisfies them too, induction on stage gives equality
on the bases and then on each constructor value; hence \(h'=h\).
\end{proof}
\end{annexitem}

\begin{annexitem}{annex:p1.10}{thm:p01-rtc-path}{Reflexive-transitive closure paths}
\begin{proof}
Let \(S\) relate \(x\) to \(y\) iff a finite \(R\)-path of the stated form
exists.  A length-zero path makes \(S\) reflexive; a length-one path shows
\(R\subseteq S\); and concatenating two paths at their common endpoint makes
\(S\) transitive.  Since \(R^*\) is the least reflexive transitive relation
containing \(R\), \(R^*\subseteq S\).

Conversely, induction on the path length proves \(S\subseteq R^*\).  A
length-zero path has equal endpoints, so reflexivity applies.  For a path of
length \(n+1\), the initial segment gives \(xR^*x_n\) by induction, while
\(x_nRx_{n+1}\) and \(R\subseteq R^*\); transitivity gives
\(xR^*x_{n+1}\).  Hence \(S=R^*\).
\end{proof}
\end{annexitem}
